Prescribed-measure existence does not by itself bound the difference from the reference. A superlinear capacity bound on the right-hand measure controls that difference. The sublevel estimate from \cref{prop:relative-capacity-sublevel-estimate} reduces the argument to an iteration with summable step sizes.

\begin{theorem}\label{thm:relative-capacity-dominated-boundedness}
Suppose $u\in\mathcal E(X,\theta,\phi)$ and $u\leq\phi$. Let $0\leq a<1$, $A>0$, $\delta>0$, and let $\nu$ be a finite positive measure such that
\begin{equation}\label{eq:relative-capacity-measure-domination}
 \mu_u\leq\nu+a\mu_\phi,\qquad
 \nu(E)\leq A\operatorname{Cap}_{\theta,\phi}(E)^{1+\delta}
                    \quad\text{for every Borel }E\subset X.
\end{equation}
Then $\phi-C\leq u\leq\phi$ for some finite $C$. If $I=\int_X(\phi-u)\,\mathrm d\nu<\infty$, one can take
\begin{equation}\label{eq:relative-capacity-moment-depth}
 C=2+\frac I{1-a}\left(2\left(\frac A{1-a}\right)^{1/n}\right)^{n/\delta}
                         +\frac1{1-2^{-\delta}}.
\end{equation}
Differences in these integrals and sublevel sets are interpreted off the common polar exceptional sets.
\end{theorem}
\begin{proof}
Write $E_s=\{u<\phi-s\}$, $g(s)=\operatorname{Cap}_{\theta,\phi}(E_s)^{1/n}$ and $B=(A/(1-a))^{1/n}$. Relative full-mass comparison in \cref{prop:relative-mixed-comparison}, applied to $u,\phi-s$, gives $\mu_\phi(E_s)\leq\mu_u(E_s)$. Hence \eqref{eq:relative-capacity-measure-domination} implies
\[
 (1-a)\mu_u(E_s)\leq\nu(E_s).
\]
For $0<t\leq1$, the earlier sublevel estimate now gives
\begin{equation}\label{eq:relative-capacity-iteration}
 t g(s+t)\leq B g(s)^{1+\delta},\qquad s>0.
\end{equation}
That same estimate with $t=1$ shows $g(s)\to0$ as $s\to\infty$, since $u$ has relative full mass. Choose $s_0>0$ with $2Bg(s_0)^\delta\leq1$, and put
\begin{equation}\label{eq:relative-capacity-explicit-depth}
 C=s_0+\frac{2B}{1-2^{-\delta}}g(s_0)^\delta.
\end{equation}

If $g(s_0)=0$, set $C=s_0$. Otherwise, as long as $g(s_j)>0$, define
\[
 t_j=2B g(s_j)^\delta,\qquad s_{j+1}=s_j+t_j.
\]
Inductively $t_j\leq1$, and \eqref{eq:relative-capacity-iteration} gives $g(s_{j+1})\leq g(s_j)/2$. If this process stops, its final level has zero capacity. If it continues, the levels increase to a finite $S$, with
\[
 S\leq s_0+2B g(s_0)^\delta\sum_{j=0}^\infty2^{-j\delta}=C.
\]
Monotonicity gives $g(S)\leq g(s_j)\to0$ and thus $g(C)=0$ in either case. The capacity bound makes $\nu(E_C)=0$; since $\phi$ itself is a strip test, $\mu_\phi(E_C)=0$ too. The measure inequality gives $\mu_u(E_C)=0$. Apply \cref{thm:relative-full-mass-domination} to $u,\phi-C$ to obtain the pointwise bound $u\geq\phi-C$.

When $I<\infty$, Markov's inequality and the sublevel estimate give
\[
 g(s+1)^n\leq\frac{\nu(E_s)}{1-a}
                          \leq\frac I{(1-a)s}.
\]
Choose $s=1+I(2B)^{n/\delta}/(1-a)$ and $s_0=s+1$. Then $2Bg(s_0)^\delta\leq1$, so the last term in \eqref{eq:relative-capacity-explicit-depth} is at most $(1-2^{-\delta})^{-1}$. This proves \eqref{eq:relative-capacity-moment-depth}.
\end{proof}
\index{Monge--Amp\`ere equation!capacity boundedness estimate}

\Needspace{14\baselineskip}
\begin{corollary}\label{cor:relative-capacity-energy-boundedness}
Suppose $u\in\mathcal E^1(X,\theta,\phi)$, $u\leq\phi$, $A>0$, $\delta>0$, and
\[
 \mu_u(E)\leq A\operatorname{Cap}_{\theta,\phi}(E)^{1+\delta}
                         \quad\text{for every Borel }E\subset X.
\]
If $L\geq0$ and $-E_\theta^\phi(u)\leq L$, then
\[
 \phi-\left(2+(n+1)L(2A^{1/n})^{n/\delta}
                        +\frac1{1-2^{-\delta}}\right)\leq u\leq\phi.
\]
\end{corollary}
\begin{proof}
Use the theorem with $a=0$, $\nu=\mu_u$, and the bound
$\int_X(\phi-u)\mu_u\leq(n+1)L$ from \eqref{eq:relative-energy-integrability-bounds}.
\end{proof}

\begin{corollary}\label{cor:relative-lp-capacity-boundedness}
Let $p>1$, $0\leq f\in L^p(X,\omega^n)$, and $u\in\mathcal E(X,\theta,\phi)$ with $\sup_Xu=0$. Suppose $0\leq a<1$ and
\[
 \mu_u\leq f\omega^n+a\mu_\phi.
\]
Then $\phi-C\leq u\leq\phi$, with $C$ depending only on $X,\theta,\omega,p,m,\|f\|_{L^p}$ and $(1-a)^{-1}$. In particular, when $\int_X f\omega^n=m$, there is a normalized solution of $\mu_u=f\omega^n$ with bounded difference from $\phi$.
\end{corollary}
\begin{proof}
Model normalization gives $u\leq\phi$ by \cref{lem:model-normalized-supremum}; both suprema are zero. Let $q=p/(p-1)$. The uniform bound of \cref{cor:normalized-global-psh-lp-bound} and H\"older's inequality give
\[
 I=\int_X(\phi-u)f\omega^n
       \leq\|f\|_{L^p}\|\phi-u\|_{L^q}
       \leq2M_q\|f\|_{L^p}.
\]
By \cref{cor:relative-lp-capacity-domination} with $r=2$, the measure $\nu=f\omega^n$ satisfies \eqref{eq:relative-capacity-measure-domination} with $\delta=1$ and $A=B_2\|f\|_{L^p}/m^2$; enlarge $A$ to be positive if necessary. Apply \eqref{eq:relative-capacity-moment-depth}. Its constant has precisely the stated dependence. The measure $f\omega^n$ does not charge pluripolar sets, so \cref{thm:nonpluripolar-relative-prescribed-measure-existence} supplies a normalized solution when its mass is $m$. The bound just proved gives the asserted bounded difference from $\phi$.
\end{proof}
